% s04 The two cube conventions: interlacing and margins
\section{The two cube conventions: interlacing and margins}

King's printed state lives on King's convention $\NDK_{k,J}$ (\cref{def:s02:Kcube}); the rotation part and the Wick-tube estimates of this paper use the comparison convention $\NDa_{k,J}$ (\cref{def:s02:acube}). We show that the two families are interlaced in $J$ (\cref{lem:s04:G}), draw the consequences for the limits (\cref{cor:s04:G1,cor:s04:G2}), and show that the comparisons of a cube with a reflected cube, a larger box and a tiling by cells carry over to King's convention with the printed margins (\cref{lem:s04:C,lem:s04:D,lem:s04:E}).

\subsection{The two conventions and the support condition}

Throughout, $\eps=\eps_k=5^{-k}$, $\XJ=\CJ\cap\Lk=\{-MJ,\dots,MJ-\eps\}^3$, $\TJ=\{x\in\XJ:\ x_\mu=MJ-\eps\text{ for some }\mu\}$ and $\alive^K_J=\XJ\setminus\TJ=\{-MJ,\dots,MJ-2\eps\}^3$. Both states carry the same photon variables $A_\alpha(x)$, $x\in\XJ$, with the Neumann form of $\XJ$, so their photon Gaussians are identical and every plaquette observable is the same function of the shared variables. They differ exactly by the matter on $\TJ$: alive in the comparison convention (scalar Dirichlet planes $x_\mu=-MJ-\eps$, $x_\mu=MJ$), killed in King's (planes $x_\mu=-MJ-\eps$, $x_\mu=MJ-\eps$, $2MJ$ apart, a multiple of every block scale; cf.\ \cref{rem:s03:cube-identification}). A bond with one end alive and the other killed or outside the site set contributes its diagonal part $\frac\eps2r_x^2$ at the alive end, so an alive site carries six diagonals in either convention. In the proof of his Theorem~3.2 and in his box form (3.38), King \cite{K2} works in the comparison convention: every site of the region is alive there.

For $(f,g)$ with $f$ a real $2$-form and $g\ge0$ a scalar, both in $C^\infty_c$, we write as in \cref{def:s02:sources} $\Theta(\mathrm if,g)=e^{\mathrm iF(f)+W(g)}$ (class $\cI$), $\Theta(f,-g)=e^{F(f)-W(g)}$ (class $\cD$), $W(g)=\sum_x\eps^3g(x)(|\phi(x)|^2-2\sW_\eps)$ (free Wick constant), and $Z^a_{k,J}(\mathrm if,g)=e^{-a_k\sS(g)}\av{\Theta(\mathrm if,g)}$, $Z^a_{k,J}(f,-g)=e^{a_k\sS(g)}\av{\Theta(f,-g)}$, $\sS(g)=\sum_x\eps^3g(x)$; a superscript $K$ ($Z^{a,K}_{k,J}$) marks King's convention. All these expectations are real. The sources lie inside $\XJ$ if
\begin{equation}\label{eq:s04:S}\tag{S}
\begin{gathered}
\text{every plaquette $p=(x;\alpha,\beta)$ of $\Lk$ with $f_{c(p)}(e_\alpha,e_\beta)\ne0$ has its four corners in $\XJ$, and}\\
\text{every site $x\in\Lk$ with $g(x)\ne0$ lies in $\alive^K_J=\XJ\setminus\TJ$.}
\end{gathered}
\end{equation}
If $\supp f\cup\supp g\subset\{y:\dist_\infty(y,\R^3\setminus\CJ)>1\}$, as every hypothesis of this paper demands of its sources, then \eqref{eq:s04:S} holds: a site with $g(x)\ne0$ has $-MJ+1<x_\mu<MJ-1\le MJ-\eps$, so $x\in\alive^K_J$; a plaquette with $f_{c(p)}\ne0$ has its centre at $\ell^\infty$-distance $>1$ from $\R^3\setminus\CJ$ and its corners within $\eps/2\le\frac12$ of it, hence in $\XJ$. Also \eqref{eq:s04:S} for $X_{J_0}$ implies \eqref{eq:s04:S} for every $J\ge J_0$.

\subsection{Cut-outs and the interlacing lemma}

\begin{lemma}[cut-outs from a box with killed sites]\label{lem:s04:G3}
Let $\av{\cdot}_Y$ be the state of \cref{def:s02:acube} on a finite set $V\subset\Lk$ of sites with the matter killed on $T\subset V$, and $X\subset V$ any subset. Let $(f,g)$, $g\ge0$, satisfy \textup{(S$_X$)}: every plaquette with $f\ne0$ has its four corners in $X$ and every site with $g\ne0$ lies in $X$ ($g$ on $X\cap T$ is immaterial). Then deleting from $\av{\cdot}_Y$ the photon gradient terms coupling $X$ to $V\setminus X$ and the cosine parts of the hopping terms between $X$ and $V\setminus X$ yields a product whose $X$-factor is $\av{\cdot}^{(a),X\cap T}_X$, the state of \cref{def:s02:acube} on $X$ with the matter killed on $X\cap T$; consequently
\[
\av{\Theta(\mathrm if,g)}^{(a),X\cap T}_X\ \le\ \av{\Theta(\mathrm if,g)}_Y,\qquad
\av{\Theta(f,-g)}^{(a),X\cap T}_X\ \ge\ \av{\Theta(f,-g)}_Y,
\]
and the same for $Z^a$. In particular, \textup{(a)} if $X\cap T=\emptyset$ the factor is $\av{\cdot}^{(a)}_X$; \textup{(b)} if $X=\prod_\mu\{a_\mu,\dots,b_\mu\}$ is a box and $X\cap T=\{x\in X:\ x_\mu=b_\mu\text{ for some }\mu\}$, the factor is the state of King's convention on $X$.
\end{lemma}

\begin{proof}
Each deletion is in the cone direction: deleting a photon gradient term does not increase $\av{\Theta_{F,g}}$ and does not decrease $\av{\Theta^-_{F,g}}$ (\cref{lem:s02:BFS}(a),(d)), and setting a hopping coefficient $c_{xy}$ to $0$ likewise (\cref{lem:s02:BFS}(b),(d)); every intermediate measure is a BFS measure. After all deletions nothing couples the $X$-variables to the rest, so the measure is a product (\cref{lem:s02:BFS}(e)). In the $X$-factor the photon is the Neumann photon of $X$; a site of $X\cap T$ carries its point mass; an alive $x\in X\setminus T$ carries the single-site weight times $\frac\eps2r_x^2$ for each of its six bonds, since in $\av{\cdot}_Y$ every bond at an alive $x$ contributes $\frac\eps2r_x^2$ at $x$ whether its other end is alive in $V$ (the diagonal part of $\frac\eps2|e^{\mathrm ie\eps A}\phi_y-\phi_x|^2$), killed, or outside $V$ ($\phi=0$ there), and the cut keeps all diagonals; the cosines are those of the bonds within $X$ between alive sites. That is the state of \cref{def:s02:acube} on $X$ killed on $X\cap T$. By (S$_X$) the observable is a function of the $X$-factor's variables (at $x\in X\cap T$ the factor $e^{\eps^3g(x)r_x^2}$ is $1$ in both states), and the $c$-numbers $e^{\mp2\sW_\eps\sS(g)}$, $e^{\mp a_k\sS(g)}$ are the same positive constants on both sides (\cref{lem:s02:alg}(a); the counterterm is site-independent). Hence $\av{\Theta}^{(a),X\cap T}_X$ is the expectation of $\Theta$ in the fully deleted measure, which is $\le\av{\Theta}_Y$ on $\cI$ and $\ge$ on $\cD$.
\end{proof}

The clause of \cref{lem:s02:cut} that $X$ and its neighbours lie in one fundamental cell is a torus clause; for a box it has no content, since a bond leaving $V$ already carries its Dirichlet diagonal in $\av{\cdot}_Y$. Every sandwich $C_{J-1}\subset\cdot\subset C_{J+1}$ of this paper, read in King's convention, is an instance of \cref{lem:s04:G3}.

\begin{lemma}[interlacing of the two cube families]\label{lem:s04:G}
Fix $k\ge0$, $J\ge1$ and $(f,g)$, $g\ge0$, satisfying \eqref{eq:s04:S} for $\XJ$. Then
\[
\avK{\Theta(\mathrm if,g)}\ \le\ \ava{\Theta(\mathrm if,g)}\ \le\ \langle\Theta(\mathrm if,g)\rangle^{K}_{k,J+1},\qquad
\avK{\Theta(f,-g)}\ \ge\ \ava{\Theta(f,-g)}\ \ge\ \langle\Theta(f,-g)\rangle^{K}_{k,J+1},
\]
and the same for the scheme functionals $Z^a_{k,J}(\mathrm if,g)$, $Z^a_{k,J}(f,-g)$ of the three states.
\end{lemma}

\begin{proof}
\emph{Step 0 (one observable, one multiplier).} By \eqref{eq:s04:S} the plaquettes with non-zero coefficient have all corners in $\XJ$, whose photon variables are present in all three states, so $F(f)$ is one real linear functional $F$ of the shared variables; the sites with $g\ne0$ lie in $\alive^K_J$, alive in all three states; $\sW_\eps$, $\dm_\eps$, $\sS(g)$ are the same for the three regions. Hence, with $g'_x:=\eps^3g(x)\ge0$, $\Theta(\mathrm if,g)=e^{-2\sW_\eps\sS(g)}\Theta_{F,g'}$ and $\Theta(f,-g)=e^{2\sW_\eps\sS(g)}\Theta^-_{F,g'}$, where $\Theta_{F,g'}=e^{\mathrm iF}e^{\sum_xg'_xr_x^2}$, $\Theta^-_{F,g'}=e^{F}e^{-\sum_xg'_xr_x^2}$ are the observables of \cref{lem:s02:BFS}, and the scheme multiplier $e^{\mp a_k\sS(g)}$ is one positive constant for all three states (\cref{lem:s02:alg}(a)). It suffices to treat $\av{\Theta_{F,g'}}$, $\av{\Theta^-_{F,g'}}$.

\emph{Step 1 (kill step on $\cI$).} For $s\ge0$ let $\av{\cdot}_s$ be $\ava{\cdot}$ with the single-site measure at each $x\in\TJ$ multiplied by $e^{-s\eps^3r_x^2}$, a BFS measure, and $\varphi(s):=\av{\Theta_{F,g'}}_s=N(s)/D(s)$, $N(s)=\int\Theta\,e^{-s\eps^3\sum_{\TJ}r_x^2}d\mu_0$, $D(s)=\int e^{-s\eps^3\sum_{\TJ}r_x^2}d\mu_0$, $\mu_0$ the unnormalised measure at $s=0$. The $s$-derivative of either integrand is bounded in modulus by $\eps^3\sum_{\TJ}r_x^2$ times the integrand at $s=0$, integrable against the quartic weight $e^{-\lambda\eps^3r^4}$ and the Gaussian; so differentiation under the integral sign is justified, $D(s)>0$, and
\[
\varphi'(s)=-\eps^3\sum_{x\in\TJ}\big(\av{\Theta\,r_x^2}_s-\av{\Theta}_s\av{r_x^2}_s\big)=-\eps^3\sum_{x\in\TJ}\av{\Theta;r_x^2}_s\ \le\ 0
\]
by \cref{lem:s02:BFS}(c). So $\varphi(s)\le\varphi(0)=\ava{\Theta}$. Since $\Theta_{F,g'}$ is continuous with $|\Theta_{F,g'}|=e^{\sum_xg'_xr_x^2}$ and $\ava{\cdot}$ is of King type, \cref{lem:s02:Kprime} gives $\varphi(s)\to\av{\Theta}_{\TJ}$, the state with the matter killed on $\TJ$, which is $\avK{\cdot}$. Hence $\avK{\Theta}\le\ava{\Theta}$.

\emph{Step 2 (cut step).} Apply \cref{lem:s04:G3} with $V:=X_{J+1}$, $T:=T_{J+1}$ (so $\av{\cdot}_Y=\langle\cdot\rangle^K_{k,J+1}$), $X:=\XJ$: $\XJ\subset X_{J+1}$, and $\XJ\cap T_{J+1}=\emptyset$ since every site of $T_{J+1}$ has a coordinate $M(J+1)-\eps\ge MJ>MJ-\eps$. (At $k=0$, $M=1$ the outward neighbours of $\XJ$ lie on $T_{J+1}$ and are killed in $\av{\cdot}_Y$; harmless, since the bond from an alive $x\in\XJ$ to such a neighbour already carries $\frac\eps2r_x^2$ at $x$, the Dirichlet diagonal of the comparison cube.) By \cref{lem:s04:G3}(a) the $\XJ$-factor is $\ava{\cdot}$, and \eqref{eq:s04:S} is (S$_{\XJ}$): $\ava{\Theta(\mathrm if,g)}\le\langle\Theta(\mathrm if,g)\rangle^K_{k,J+1}$, $\ava{\Theta(f,-g)}\ge\langle\Theta(f,-g)\rangle^K_{k,J+1}$.

\emph{Step 3 (kill step on $\cD$).} With $\av{\cdot}_s$ as in Step~1 and $\varphi^-(s):=\av{\Theta^-_{F,g'}}_s$ (differentiation justified as before, now with $|\Theta^-|\le e^{F(A)}$ integrable against the Gaussian), $(\varphi^-)'(s)=-\eps^3\sum_{x\in\TJ}\av{\Theta^-;r_x^2}_s\ge0$ by \cref{lem:s02:BFS}(d), so $\varphi^-(s)\ge\varphi^-(0)=\ava{\Theta^-}$; as $|\Theta^-_{F,g'}|\le e^{F(A)}$ with $F$ real linear, \cref{lem:s02:Kprime} (with $\ell=F$, $C=0$) gives $\varphi^-(s)\to\avK{\Theta^-}$. Hence $\avK{\Theta(f,-g)}\ge\ava{\Theta(f,-g)}$.

\emph{Step 4 ($Z^a$).} The multiplier of Step~0 is one positive constant.
\end{proof}

Nothing in the proof uses a reflection symmetry of either cube (neither has one), the counterterm beyond its site-independence, or the renormalisation group. Since the observable vanishes at the killed sites in every kill step of this section, each such step also follows with no limit: set $c_{xy}=0$ on every bond at the set $U$ to be killed, one bond at a time (\cref{lem:s02:BFS}(b),(d)); the measure is then a product (\cref{lem:s02:BFS}(e)) whose $U$-factor cancels, and the other factor is the state with $U$ killed.

\subsection{Monotonicity of the K family and the limits}

\begin{corollary}[the K family is monotone in $J$]\label{cor:s04:G1}
Under \eqref{eq:s04:S} for $\XJ$: $\avK{\Theta(\mathrm if,g)}\le\langle\Theta(\mathrm if,g)\rangle^K_{k,J+1}$ and $\avK{\Theta(f,-g)}\ge\langle\Theta(f,-g)\rangle^K_{k,J+1}$, and the same for $Z^{a,K}$; hence for all $J\ge J_0$ whenever the sources satisfy \eqref{eq:s04:S} for $X_{J_0}$.
\end{corollary}

\begin{proof}
Chain the two inequalities of \cref{lem:s04:G}. (For the comparison family this is \cref{lem:s04:G3}(a) with $V=X_{J+1}$, $X=\XJ$, \ie the first inequality of \cref{lem:s02:cut}.)
\end{proof}

The proof of Theorem~3.2 of \cite{K2} first replaces the $J$-dependent Dirichlet Wick subtraction (2.13) by the free, $J$-independent one and then proves the monotonicity in $J$ in the comparison convention (\cref{lem:s04:G3}(a)); by \cref{cor:s04:G1} that statement holds between cubes of King's convention too, with the kill step (Step~1) not printed there. The form with the Dirichlet constant follows in neither convention and is not used here; at $g=0$ no Wick constant enters.

\begin{corollary}[the limits]\label{cor:s04:G2}
Let $f=(f_1,\dots,f_N)$ be real $2$-forms in $C^\infty_c$ with $\supp f_j\subset\{\dist_\infty(\cdot,\R^3\setminus C_{J_0})>1\}$; for $J\ge J_0$, $z\in\C^N$ put $a^K_{k,J}(z):=\avK{e^{\sum_jz_jF(f_j)}}$, $a^{(a)}_{k,J}(z):=\ava{e^{\sum_jz_jF(f_j)}}$, $g(z):=\sum_j(\Re z_j)f_j$. At the cone points $z=\mathrm it$, $t\in\R^N$, both are real with modulus $\le1$, and \eqref{eq:s04:S} holds for $\XJ$.
\begin{enumerate}[label=\textup{(\roman*)},leftmargin=*]
\item \textup{(One $J$-limit at fixed $k$; no hypothesis.)} For every $k$ and $t$, $J\mapsto a^K_{k,J}(\mathrm it)$ and $J\mapsto a^{(a)}_{k,J}(\mathrm it)$ are nondecreasing on $J\ge J_0$, bounded by $1$ and interlaced; hence $\lim_Ja^K_{k,J}(\mathrm it)=\lim_Ja^{(a)}_{k,J}(\mathrm it)\in[-1,1]$. The same holds for $\av{\Theta(\mathrm if,g)}$, $g\ge0$, with a common limit in $(-\infty,\infty]$.
\item \textup{(King's cube limits.)} By \cref{hyp:s03:EC}, posed on King's convention, for every such family, every $J\ge J_0$ and every $t\in\R^N$, $\lim_ka^K_{k,J}(t)$ exists in $(0,\infty)$. Then: ($\alpha$) $a^K_{k,J}\to\ZKJ$ locally uniformly on $\C^N$, $\ZKJ$ is entire and $|\ZKJ(z)|\le e^{\norm{g(z)}_2^2}$; ($\beta$) $\ZKJ\to\ZKinf$ locally uniformly on $\C^N$, $\ZKinf$ is entire and $|\ZKinf(z)|\le e^{\norm{g(z)}_2^2}$; ($\gamma$) for all $J\ge J_0$ and $t\in\R^N$,
\[
\ZKJ(\mathrm it)\ \le\ \liminf_ka^{(a)}_{k,J}(\mathrm it)\ \le\ \limsup_ka^{(a)}_{k,J}(\mathrm it)\ \le\ Z^K_{J+1}(\mathrm it),
\]
hence $\lim_J\liminf_ka^{(a)}_{k,J}(\mathrm it)=\lim_J\limsup_ka^{(a)}_{k,J}(\mathrm it)=\ZKinf(\mathrm it)$.
\item \textup{(The subsequential limits of the comparison family.)} Let $\cS=(k_m)$ be an increasing sequence of spacings along which $a^{(a)}_{k_m,J}\to Z^{\cS}_J$ locally uniformly on $\C^N$ for every $J\ge J_0$ and every such family (a continuum subsequence of the comparison family, \cref{def:s12:contsub}; they exist with no hypothesis, \cref{lem:s12:LIM0}). Then $Z^{\cS}_J\to\ZKinf$ locally uniformly on $\C^N$ as $J\to\infty$: every subsequential fixed-cube limit of the comparison family, followed by $J\to\infty$, is King's printed state. In particular, if the fixed-cube limits of the comparison family exist along the full sequence, $\lim_J\lim_ka^{(a)}_{k,J}=\ZKinf$.
\end{enumerate}
\end{corollary}

\begin{proof}
(i) Monotonicity is \cref{cor:s04:G1} for the K family and \cref{lem:s04:G3}(a) for the comparison family; \cref{lem:s04:G} interlaces them; two nondecreasing sequences with $a_J\le b_J\le a_{J+1}$ have the same limit, finite or $+\infty$.

(ii)($\alpha$) This is the first step of the proof of \cref{thm:s05:C}(b) run on the K family. (1) $a^K_{k,J}$ is entire (Gaussian decay in $A$, $e^{-\lambda\eps^3r^4}$ decay in $r$ at the alive sites). (2) The local bound of \cref{lem:s05:C2} holds for the K family with the same constant: $\Re\sum_jz_jF(f_j)=F(g(z))$, so $|a^K_{k,J}(z)|\le\avK{e^{F(g(z))}}$; \cref{lem:s02:GD}(a) with $s=1$ and $v:=d^*_kg(z)$ (defined by $F_{\Lk}(g(z))=\ip{v}{A}_k$) applies to any BFS measure whose Gaussian has covariance $C$ and gives $\av{e^{\ip{v}{A}}}\le e^{\ip{v}{Cv}/2}$; the K cube is a BFS measure whose Gaussian is the Neumann photon Gaussian of $\XJ$, the same as the comparison cube's, so its covariance is the $C$ of \cref{lem:s05:C1}, a statement about this $C$ and the plaquettes with corners in $\XJ$: $\ip{v}{Cv}_k\le2\norm{g(z)}_{2,k}^2$. Hence $|a^K_{k,J}(z)|\le e^{\norm{g(z)}_{2,k}^2}\le e^{3C_0\fournorm{g(z)}^2}$ uniformly in $k\ge0$, $J\ge J_0$, by (E3) of \cref{lem:s02:E}: the plaquette centres of each of the three orientations form a translate of $\eps_k\Z^3$, each contributing at most $C_0\fournorm{g}^2$ to $\norm{g}_{2,k}^2$ by (E1), and the three orientations give $3C_0$. Nothing in this chain sees the matter beyond the state being a BFS measure. (3) \Cref{hyp:s03:EC} gives convergence at every real point. (4) $\R^N$ is a set of uniqueness for entire functions on $\C^N$; by Montel's theorem every subsequence has a locally uniformly convergent sub-subsequence, all limits are entire and agree on $\R^N$, hence everywhere \cite{Hor90}. So $a^K_{k,J}\to\ZKJ$ locally uniformly, $\ZKJ$ is entire, and the bound passes to the limit, $|\ZKJ(z)|\le\lim_ke^{\norm{g(z)}_{2,k}^2}=e^{\norm{g(z)}_2^2}$ (\cref{lem:s02:E}), a $J$-free bound.

($\beta$) By \cref{cor:s04:G1} at every $k$ and ($\alpha$), $J\mapsto\ZKJ(\mathrm it)$ is nondecreasing and $\le1$, so $\ZKinf(\mathrm it):=\lim_J\ZKJ(\mathrm it)$ exists; $\{\ZKJ\}_{J\ge J_0}$ is locally uniformly bounded by the $J$-free bound of ($\alpha$); $(\mathrm i\R)^N$ is a maximal totally real subspace, hence a set of uniqueness for entire functions; Montel's theorem and the identity theorem give $\ZKJ\to\ZKinf$ locally uniformly, $\ZKinf$ entire, $|\ZKinf(z)|\le e^{\norm{g(z)}_2^2}$. (This is the proof of \cref{thm:s05:C}(b) with \cref{lem:s02:cut} replaced by \cref{cor:s04:G1}.)

($\gamma$) \cref{lem:s04:G} at every $k$, with $k\to\infty$ at the two ends (limits by ($\alpha$)) and $\liminf$, $\limsup$ in the middle, gives the sandwich; ($\beta$) gives the limits.

(iii) \cref{lem:s04:G} along $\cS$, with the full-sequence limits of ($\alpha$) at the ends, gives $\ZKJ(\mathrm it)\le Z^{\cS}_J(\mathrm it)\le Z^K_{J+1}(\mathrm it)$, so $\lim_JZ^{\cS}_J(\mathrm it)=\ZKinf(\mathrm it)$ by ($\beta$). $\{Z^{\cS}_J\}_J$ is locally bounded by the $J$-free bound $e^{\norm{g(z)}_2^2}$ (\cref{lem:s05:C2} passed to the limit), so by Montel's theorem and the uniqueness set $(\mathrm i\R)^N$ it converges locally uniformly to an entire function (\cref{lem:s12:LIM1}), which agrees with $\ZKinf$ on $(\mathrm i\R)^N$, hence on $\C^N$. If the fixed-cube limits exist along the full sequence, that sequence is a continuum subsequence.
\end{proof}

Statements (ii)--(iii) also hold with the two conventions exchanged, for $J\ge J_0+1$, when the comparison family is convergent at fixed cubes (\cref{def:s03:conv}; read \cref{lem:s04:G} as $a^{(a)}_{k,J-1}\le a^K_{k,J}\le a^{(a)}_{k,J}$); the exchanged (iii) needs \cref{lem:s12:LIM0,lem:s12:LIM1} for the K family, whose proofs use only the local bound and the photon-only bound $\av{F(u-u')^2}\le2\norm{u-u'}_{2,k}^2$, both transferred by the covariance identity in ($\alpha$).

\emph{What does not follow.} At fixed $J$, \cref{lem:s04:G} yields exactly ($\gamma$) and nothing on whether $\lim_ka^{(a)}_{k,J}(\mathrm it)$ exists: the bounds $\ZKJ(\mathrm it)$, $Z^K_{J+1}(\mathrm it)$ differ in general (the cubes differ by $M$, so the gap does not shrink with $k$), and nothing constrains the oscillation of $a^{(a)}_{k,J}(\mathrm it)$ in $k$ between them. So the existence of the fixed-cube limits on one convention does not imply it on the other by these inputs, and likewise with scalar sources: the existence theorem must be proved on the family on which the renormalisation-group construction is carried out; the other is a comparison family with the same infinite-volume state by (ii)--(iii). No counterexample is claimed.

\subsection{Three margins in King's convention}

Let $\sigma_3=\diag(1,1,-1)$, $X:=\XJ$, $X_+:=\sigma_3X=(\CJ+\eps e_3)\cap\Lk$ (third coordinates $\{-MJ+\eps,\dots,MJ\}$), $\sigma_3\TJ=\{x\in X_+:\ x_1=MJ-\eps\text{ or }x_2=MJ-\eps\text{ or }x_3=-MJ+\eps\}$ and $A_\sigma:=\sigma_3(\alive^K_J)=\{x:\ x_1,x_2\in\{-MJ,\dots,MJ-2\eps\},\ x_3\in\{-MJ+2\eps,\dots,MJ\}\}$.

\begin{lemma}[the reflection sandwich]\label{lem:s04:C}
Let $(f,g)\in\cK$ be supported at $\ell^\infty$-distance $>1$ from $\R^3\setminus C_{J_0}$.
\textup{(a)} If $M\ge2\eps_k$ (\ie $k\ge1$, or $M\ge2$), then for every $J\ge J_0+1$
\[
Z^{a,K}_{k,J-1}(\mathrm if,g)\ \le\ Z^{a,K}_{k,J}(\mathrm i\,\sigma_3\!\cdot\!f,\sigma_3\!\cdot\!g)\ \le\ Z^{a,K}_{k,J+1}(\mathrm if,g).
\]
\textup{(b)} For every $k\ge0$, $M\ge1$ and $J\ge J_0+2$ the same holds with $J\mp2$ in place of $J\mp1$.
\textup{(c)} For $k=0$, $M=1$ the one-step sandwich \textup{(a)} is not obtainable by deletions and kills: $A_\sigma$ contains alive sites ($x_3=MJ$) killed in King's cube of index $J+1$, and $\alive^K_{J-1}$ contains sites ($x_3=-MJ+\eps$) killed in the reflected state. Whether the inequality itself holds there is left open.
\end{lemma}

\begin{proof}
\emph{Step 1 (transport).} By \cref{lem:s06:PK}, $Z^{a,K}_{k,J}(\mathrm i\sigma_3\!\cdot\!f,\sigma_3\!\cdot\!g)=e^{-a_k\sS(g)}\av{\Theta(\mathrm if,g)}_R$, with $R$ the state on $X_+$ for the frame $(e_1,e_2,-e_3)$ killed on $\sigma_3\TJ$ ($\sS(\sigma_3\cdot g)=\sS(g)$). In the frame $(e_\alpha)$, as in \cref{lem:s02:refl}, $R$ is the BFS measure with $A_1,A_2$ on $X_+$ and $A_3$ on $X_+-\eps e_3=X$ (gradient pairs within these sets, masses on them), matter alive exactly on $A_\sigma$ with point masses on $\sigma_3\TJ$ and $\phi=0$ outside $X_+$, all six diagonals at each alive site, and the standard cosine on every bond with both ends in $A_\sigma$ (the flipped bond with $A'_3(x)=-A_3(x-\eps e_3)$ gives the standard cosine of $\langle x-\eps e_3,x\rangle$); $F(f)$, $W(g)$ are the standard ones (the plaquettes $(x;\alpha,3)$ pick up two sign changes and keep their centres). $R$ is of King type. The sources satisfy \eqref{eq:s04:S} for $X_{J-1}$ (resp.\ $X_{J-2}$), so by the inclusions below Step~0 of \cref{lem:s04:G} applies to all states compared.

\emph{Step 2 (upper end).} If $A_\sigma\subset\alive^K_{J'}$ and $X_+\cup X\subset X_{J'}$: from $\langle\cdot\rangle^K_{k,J'}$ kill every site of $\alive^K_{J'}\setminus A_\sigma$ (\cref{lem:s02:BFS}(c), \cref{lem:s02:Kprime}, which allows the already killed $T_{J'}$), leaving exactly the cosines of bonds within $A_\sigma$, \ie $R$'s; then delete every photon gradient term except the pairs within $X_+$ for $A_1,A_2$ and within $X$ for $A_3$ (\cref{lem:s02:BFS}(a); these are gradient pairs of $X_{J'}$). The result is a product (\cref{lem:s02:BFS}(e)) whose factor on $R$'s variables is $R$ (photon, alive set with all six diagonals by the bookkeeping of \cref{lem:s04:G3}, cosines); the other factor cancels. Hence $\av{\Theta(\mathrm if,g)}_R\le\langle\Theta(\mathrm if,g)\rangle^K_{k,J'}$. The inclusions hold iff $MJ\le MJ'-2\eps$ (direction $3$; the rest is trivial): for $J'=J+1$ iff $M\ge2\eps$, for $J'=J+2$ always.

\emph{Step 3 (lower end).} If $\alive^K_{J''}\subset A_\sigma$ and $X_{J''}\subset X\cap X_+$: from $R$ kill every site of $A_\sigma\setminus\alive^K_{J''}$ (\cref{lem:s02:BFS}(c), \cref{lem:s02:Kprime}) and delete every photon gradient term coupling an $X_{J''}$-variable to another variable of $R$ (\cref{lem:s02:BFS}(a)). The $X_{J''}$-factor has the Neumann photon of $X_{J''}$ (gradient pairs within $X_+\cap X_{J''}=X\cap X_{J''}=X_{J''}$) and the matter alive on $\alive^K_{J''}$ with all diagonals, killed on $T_{J''}$: it is $\langle\cdot\rangle^K_{k,J''}$; the rest cancels. Hence $\langle\Theta(\mathrm if,g)\rangle^K_{k,J''}\le\av{\Theta(\mathrm if,g)}_R$. The inclusions hold iff $-MJ''\ge-MJ+2\eps$: for $J''=J-1$ iff $M\ge2\eps$, for $J''=J-2$ always.

\emph{Step 4.} Steps 2--3 with $J\pm1$ give (a), with $J\pm2$ give (b); the multipliers agree (\cref{lem:s02:alg}(a)). For (c): at $k=0$, $M=1$ the top of $A_\sigma$ in direction $3$ is $J$ while that of $\alive^K_{J+1}$ is $J-1$, and the bottom of $\alive^K_{J-1}$ is $-J+1=-MJ+\eps$, the killed layer of $R$; deletions and kills never revive a killed site.
\end{proof}

In the curvature sector ($g=0$) the limit sentences of \cref{thm:s06:S1K} run on (a) for $k\ge1$ or on (b) for all $k$, with $k\to\infty$ by \cref{cor:s04:G2}(ii) and $J\to\infty$ by \cref{cor:s04:G1}; with scalar sources they need the K-family tube limits, supplied by \cref{thm:s13:TWK} (\cref{rem:s04:wick}).

\begin{lemma}[the upper box needs no change]\label{lem:s04:D}
Let $\Omega\subset B$ be lattice boxes at spacing $\eps$ such that $B$ contains every site within $\ell^\infty$-distance $N$ lattice steps of $\Omega$, $N\ge1$ (the hypothesis of \cref{lem:s07:U}, whose mass condition forces $N\ge1$), and let $\av{\cdot}^K_B$ be the state of King's convention on $B$ (killed on its top face layers $T_B$). Then \textup{(a)} $\Omega\cap T_B=\emptyset$; \textup{(b)} the second part of \cref{lem:s07:U} holds with $\av{\cdot}^K_B$ in place of the comparison state on $B$, with the same proof, and its first part is a photon statement, unchanged; \textup{(c)} consequently the sandwich in the proof of \cref{thm:s08:OS2C} holds for King's cubes with the printed cubes and margins \eqref{eq:s07:margin}: for every $k$ and every $\zeta$ admitted there, $Z^{a,K}_{k,J}(\zeta)\le\Phi_k(\zeta)\le Z^{a,K}_{k,J+2n}(\zeta)$, $\Phi_k$ the functional of the box state of \cref{def:s07:box} on $\Omega_{J,k}$; the lower end is \cref{lem:s07:L} with Step~1 of \cref{lem:s04:G} prepended. (By \cref{cor:s04:G1} the cube of index $J+2n+1$ would also do; it is not needed.)
\end{lemma}

\begin{proof}
(a) For $x\in\Omega$ and each $\mu$, $x+N\eps e_\mu\in B$, so $x_\mu<\max_Bx_\mu$ and $x\notin T_B$. (b) The proof of \cref{lem:s07:U} starts from the state on $B$ and removes all hopping cosines touching $B\setminus\Omega$; the matter on $\Omega$ is then the Dirichlet matter of $\Omega$ (diagonals of the bonds leaving $\Omega$ kept, each such bond ending in $B$ as $N\ge1$), coupled to $A$ only through the bonds within $\Omega$, and the sites of $B\setminus\Omega$ decouple. For $\av{\cdot}^K_B$ the cosines touching $T_B$ are already absent; remove the remaining cosines touching $B\setminus\Omega$ (\cref{lem:s02:BFS}(b) on $\cI$, (d) on $\cD$). By (a) every site of $\Omega$ is alive, each bond at $x\in\Omega$ carries $\frac\eps2r_x^2$ at $x$ (its other end alive, killed or outside), and the cosines within $\Omega$ are present; so the matter on $\Omega$ is the same Dirichlet matter, and the sites of $B\setminus\Omega$, alive or killed, decouple and cancel. The rest of the proof concerns only the photon law of $B$, the same in both conventions; the first part uses only the photon forms and a cutoff supported in $B$. (c) Lower end: \cref{lem:s07:L} gives $\langle Y\rangle^{(a)}_{C_J}\le\langle Y\rangle_{\Omega_{J,k}}$ on $\cI$ (its proof lands on the Dirichlet matter of $C_J$, the comparison convention), and Step~1 of \cref{lem:s04:G} gives $\langle Y\rangle^K_{C_J}\le\langle Y\rangle^{(a)}_{C_J}$, \eqref{eq:s04:S} holding since the sources of the enlarged family and their reflections are supported at $\ell^\infty$-distance $>1$ from $\R^3\setminus C_{J_0}$, $C_{J_0}\subset C_J$. Upper end: (b) with $B=X_{J+2n}$, whose hypothesis is the second margin of \eqref{eq:s07:margin}, a statement about photon sites. The multipliers agree (\cref{lem:s02:alg}(a)).
\end{proof}

For the cell bound we use the cells of \cref{def:s11:cells}: $\ell_c=2ML_c\ge10$; the cells $c=\prod_\mu[w_\mu,w_\mu+\ell_c)$ of the tiling $\cT_s$, $s\in\{0,1\}^3$, $w_\mu\in W_\mu(s_\mu)$, $W_\mu(0)=\ell_c\Z$, $W_\mu(1)=ML_c+\ell_c\Z$ (so $w_\mu,w_\mu+\ell_c,MJ\in\Z$); the good sets $G_s=\{x:\dist(x_\mu,W_\mu(s_\mu))>2\ \forall\mu\}$ and cutoffs $\chi_s$ with $\supp\chi_s\subset G_s$; $X_c:=c\cap\XJ$; $J_1(R)=\min\{J:MJ\ge R+\ell_c\}$ for sources supported in $Q_R=[-R,R]^3$; $B_{\rm cell}$, defined on the comparison cell $C_{L_c}$. $N^K_{k,C_J}$ is the normalised scalar-source functional of $\avK{\cdot}$ and $N^{(a),T}_{k,X}$ that of the comparison state on $X$ killed on $T\subset X$.

\begin{lemma}[the cell lemmas hold as printed]\label{lem:s04:E}
\textup{(a)} For a tiling $\cT_s$ and any lattice function $q\ge0$ on $\XJ$: $N^K_{k,C_J}(0,-q)\le\prod_iN^{(a),X_i\cap\TJ}_{k,X_i}(0,-q|_{X_i})$, over the site sets $X_i=c\cap\XJ$ of the cells meeting $\CJ$.
\textup{(b)} For a full cell $c\subset\CJ$: $X_c\cap\TJ\ne\emptyset$ iff $w_\mu+\ell_c=MJ$ for some $\mu$.
\textup{(c)} Let $q\ge0$ be any lattice function supported in $Q_R$, $J\ge J_1(R)$, $s\in\{0,1\}^3$ and $c\in\cT_s$ with $(\chi_sq)|_{X_c}\not\equiv0$. Then $-MJ<w_\mu$ and $w_\mu+\ell_c<MJ-2$ for every $\mu$: $c$ is a full cell inside $\CJ$ with $X_c\cap\TJ=\emptyset$, and $(\chi_sq)|_{X_c}$ is supported at $\ell^\infty$-distance $>2$ from $\R^3\setminus c$.
\textup{(d)} Hence \cref{lem:s11:W4,lem:s11:W5,lem:s11:W67} are unchanged and \cref{prop:s11:W8} holds for $\avK{\cdot}$ as printed, under its hypotheses and with its constant $K_-$: $\ln N^K_{k,C_J}(0,-q)\le K_-\wnorm{q}$ for every $k$, every $J\ge J_1(R)$ and every lattice $0\le q\le H$ supported in $Q_R$.
\textup{(e)} (Direction.) For a cell $c$ and $\emptyset\ne T\subset X_c$: $N^{(a),T}_{k,c}(0,-\tilde q)\ge N^{(a)}_{k,c}(0,-\tilde q)$ for $\tilde q\ge0$ vanishing on $T$. So $B_{\rm cell}$, a bound on a comparison cell, does not bound a source-carrying cell of King's type, and \textup{(c)} is load-bearing.
\end{lemma}

\begin{proof}
(a) The proof of \cref{lem:s11:W3} deletes the cross-cell photon gradient terms and hopping cosines, each in the $\cD$ direction (\cref{lem:s02:BFS}(d)); by \cref{lem:s02:BFS}(e) and the factor identification of \cref{lem:s04:G3} with $V=\XJ$, $T=\TJ$, $X=X_i$, the product has the factors $N^{(a),X_i\cap\TJ}_{k,X_i}$; the observable $e^{-\sum\eps^3qr^2}$ and the multiplier $e^{a_k\sS(q)}=\prod_ie^{a_k\sS(q|_{X_i})}$ factorise, and at the sites of $\TJ$ the observable factor is $1$ in every state. (b) The sites of a full cell have $x_\mu\in\{w_\mu,\dots,w_\mu+\ell_c-\eps\}$, and $x\in\TJ$ iff $w_\mu+\ell_c-\eps=MJ-\eps$ for some $\mu$. (c) Let $x\in X_c$ with $(\chi_sq)(x)\ne0$. Then $x\in G_s$; since $w_\mu,w_\mu+\ell_c\in W_\mu(s_\mu)$ and $w_\mu\le x_\mu<w_\mu+\ell_c$, this gives $w_\mu+2<x_\mu<w_\mu+\ell_c-2$; and $q(x)\ne0$ gives $|x_\mu|\le R$. Hence $w_\mu+\ell_c<R+\ell_c-2\le MJ-2$ (as $MJ\ge R+\ell_c$) and $w_\mu>-R-\ell_c+2\ge-MJ+2$. So $c\subset(-MJ,MJ)^3$, $w_\mu+\ell_c\ne MJ$, hence $X_c\cap\TJ=\emptyset$ by (b), and $\dist_\infty(x,\R^3\setminus c)=\min_\mu\min(x_\mu-w_\mu,w_\mu+\ell_c-x_\mu)>2$. (d) In the chain of \cref{prop:s11:W8} the Hölder step of \cref{lem:s11:W67} is an inequality within one state, convention-free; then (a) with the tiling $\cT_s$ and $q_s:=8\chi_sq$. A factor with $q_s|_{X_i}\equiv0$ equals $1$ whatever is killed on $X_i$ (this covers the clipped cells at the faces of $\CJ$ and any cell flush with them); a factor with $q_s|_{X_i}\not\equiv0$ is, by (c), the comparison-cell functional $N^{(a)}_{k,c}(0,-q_s|_c)$ of a full cell $c\subset\CJ$ with $q_s|_c$ supported at distance $>2$ from $\R^3\setminus c$ and $0\le q_s\le8H=H'$: exactly the hypotheses of \cref{lem:s11:W5} (hence of \cref{lem:s11:W4}: translation to $C_{L_c}$, $B_{\rm cell}$, \cref{prop:s11:W1prime} on the comparison box), as printed; \cref{lem:s11:W67} sums the cells. (A cell flush with a face, $w_\mu+\ell_c=MJ$, can meet $\supp q$ but never $\supp(\chi_sq)$; the chain feeds the cutoff sources, not $q$, to (a).) (e) Step~3 of \cref{lem:s04:G} applied to the comparison cell with $U=T$: $s\mapsto\av{\Theta^-}_s$ is nondecreasing (\cref{lem:s02:BFS}(d)) and tends to the killed state (\cref{lem:s02:Kprime}).
\end{proof}

\begin{remark}[one-line kill steps]\label{rem:s04:F}
\Cref{prop:s11:W2} (cube $\le$ torus on $\cI$), \cref{prop:s11:W1prime} (box $\le$ cube $\le$ torus), the lower end of \cref{lem:s07:L}, and the comparison of the cube with the box state of \cref{def:s07:box} on the cone (\cref{rem:s07:K8}) are each proved for the comparison cube; for King's cube prepend Step~1 of \cref{lem:s04:G} ($\avK{\cdot}\le\ava{\cdot}$ on $\cI$), and for the $\cD$ versions (cube $\ge$ torus for negative sources) Step~3. The same applies to a box of King's type (a killed face layer) in \cref{prop:s11:W1prime}. Nothing beyond \cref{lem:s04:G} is needed.
\end{remark}

\begin{remark}[the Wick sector]\label{rem:s04:wick}
With scalar sources $g\ge0$, \cref{cor:s04:G2}(i) holds as stated. The analogues of \cref{cor:s04:G2}(ii)--(iii) on the Wick tube need the fixed-cube limits and the $(k,J)$-uniform bound \eqref{eq:s11:W} for both families. For King's family the bound is the chain of \cref{thm:s11:W}(i) run on $\NDK$, whose convention-sensitive steps are the cube-versus-torus comparison (\cref{cor:s11:W2prime}, via \cref{rem:s04:F}) and the cell bound (\cref{lem:s04:E}); the assembly of that chain and the existence of the K-family tube limits are carried out in \cref{thm:s13:TWK}, after which the tube statements hold on King's convention (\cref{rem:s11:Kconvention}, \cref{prop:s13:WK}).
\end{remark}
